% !TEX root = ../../Main.tex
\section{Survey Research: Identifying the State-of-the-Art}

The review started with surveys. A survey shows the main lines of a field faster than single papers do, and it points back to the work that first showed each result.

Hambly et al. \cite{hambly_recent_2023} gave the theoretical basis. The survey first explains reinforcement learning and neural networks on their own. It then covers how the two combine into deep reinforcement learning, and how this has been applied in finance. It gives the mathematical formulations and pseudocode for critic-only, actor-only and actor-critic algorithms, together with the papers that introduced each of them. It is the reference that was used to write \Cref{Chapter:Background}.

Fischer \cite{fischer_reinforcement_2018} had the opposite purpose. It is about practice, not theory, and it reviews implementations, not derivations. This work took from it the list of dimensions it uses to compare studies: the dataset, the network under the reinforcement learning agent, the state space, the action space, the reward function, the methods used for the bias-variance and exploration-exploitation trade-offs, and the training procedure. These dimensions became the columns of \Cref{Tables:BasicRelatedWork} and \Cref{Tables:AdvancedRelatedWork}. In this way the articles below could all be compared on the same points, and not only described one at a time.

Three more surveys were also read \cite{pricope_deep_2021, alameer_reinforcement_2022, meng_reinforcement_2019}. They agree on the general picture. Most work in the field uses equity markets, daily data and discrete trading decisions. Reported results are hard to compare, because the trading environment is different in each paper.